\documentclass[12pt]{article}
\usepackage[inner=2.0cm,outer=2.0cm,top=2.5cm,bottom=2.5cm]{geometry}
\usepackage[utf8]{inputenc}
\usepackage{amsmath,amsthm,amssymb,amsfonts,xspace}
\usepackage{longtable}
\usepackage{graphicx}
\usepackage{wrapfig}
\usepackage{mathtools}
\usepackage[export]{adjustbox}
\usepackage{mathrsfs}
\usepackage{hyperref}
\usepackage{environ}
\usepackage{xcolor}
\usepackage{listings}
\usepackage{interval}
\usepackage{enumitem}
\usepackage{parskip}
\usepackage{fancyvrb}
\usepackage{tikz}
\usepackage{multirow}
% Used for solution boxes (and aligning them)
\usepackage{mdframed}
\usepackage{changepage}
\usepackage{algorithmicx}
\usepackage{algorithm} % http://ctan.org/pkg/algorithms
\usepackage[noend]{algpseudocode} % http://ctan.org/pkg/algorithmicx
\newcommand{\homeworkTitle}{Problem Set 8: $\NP$-completeness II}
\newcommand{\homeworkDueDate}{04/13/2023 11:59pm}
\newcommand{\pitem}[1]{\item[\textbf{#1.)}]}
\newcommand{\T}{\mathrm{T}}
\newcommand{\F}{\mathrm{F}}
\newcommand{\Z}{\mathbb{Z}}
\newcommand{\N}{\mathbb{N}}
\newcommand{\R}{\mathbb{R}}
\newcommand{\Q}{\mathbb{Q}}
\renewcommand{\O}{\mathcal{O}}
\newcommand{\NP}{\mathsf{NP}}
\renewcommand{\P}{\mathsf{P}}
\newcommand{\pr}{\leq_p}
\newcommand{\Mod}[1]{\ (\text{mod}\ #1)}
\newcommand{\mmod}{\textrm{\textbf{ mod }}}
\newcommand{\powerset}{\raisebox{.15\baselineskip}{\Large\ensuremath{\wp}}
}
\newcommand{\floor}[1]{\left \lfloor #1 \right \rfloor}
\newcommand{\ceil}[1]{\left \lceil #1 \right \rceil}
\newcommand{\gauss}[3]{\begin{bmatrix} #1 \\ #2 \end{bmatrix}_{#3}}
\newcommand{\pfrac}[2]{\left(\frac{#1}{#2}\right)}
\DeclarePairedDelimiter\norm{\lVert}{\rVert}
\DeclarePairedDelimiter\abs{\lvert}{\rvert}
\newcounter{ProblemCounter}
\setcounter{ProblemCounter}{1}
\newenvironment{problem}[1][Problem]{
\begin{trivlist}
\item[\hskip 6pt \labelsep \labelsep {%
\bfseries \theProblemCounter.)%
\stepcounter{ProblemCounter}%
}]
}{
\end{trivlist} \vspace{0pt}
}
\newenvironment{solution}{
\vspace{2pt}
\begin{adjustwidth}{-3pt}{15pt}
\begin{mdframed}
\textit{Solution:} \vspace{4pt}
}{
\end{mdframed}
\end{adjustwidth}
}
\pagestyle{empty}

\begin{document}
\newpage
\noindent
\begin{center}
\framebox{
\vbox{\vspace{2mm}
\hbox to 6.28in { {\bf CS 3510: Design \& Analysis of Algorithms 
\hfill {\rm 04/05/2023}} }
\vspace{6mm}
\hbox to 6.28in { {\Large \hfill \homeworkTitle  \hfill} }
\vspace{6mm}
\hbox to 6.28in { {Prof. Abrahim Ladha \hfill 
Due: \homeworkDueDate} }
\vspace{2mm}}
}
\end{center}

\textbf{\begin{itemize}
    \item Please TYPE your solutions using Latex or any other software. Handwritten solutions \textit{won't} be accepted. 
\end{itemize}}

\begin{verbatim}
    
\end{verbatim}

\begin{center}
\underline{\textbf{Steps to write a reduction}}
\end{center}

To show that a problem is $\NP$-complete, you need to show that the problem is both in $\NP$  and $\NP$-hard. In a polynomial-time reduction, we have a \textit{Problem B}, and we want to show that it is $\NP$-complete. These are the steps:
\begin{itemize}
\item Demonstrate that \textit{problem B} is in the class $\NP$. This is a description of a procedure that verifies a candidate's solution. It needs to address the runtime. You should give a polytime verifier which takes as input a candidate problem, and a witness, and outputs true or false if the witness is a solution.
\item Demonstrate that \textit{problem B} is at least as hard as a problem previously proved to be $\NP$-complete. Choose some $\NP$-complete $A$ and prove $A \pr B$. You need to show that \textit{problem B} is $\NP$-hard. This is done via polytime reduction from a known $\NP$-complete \textit{problem A} to the unknown \textit{problem B} $(A\rightarrow B)$ as follows:
\begin{enumerate}
\item Choose $A$. You will have many $\NP$-complete problems to pick from later on, and it is best to pick a similar one. 

\item Give your reduction $f$ and argue that it runs in polynomial time.

\item Prove that $x \in A \implies f(x) \in B$

\item Prove that $x \not \in A \implies f(x) \not \in B$

\item This is sufficient to show $x \in A \iff f(x) \in B$. Since $A$ was $\NP$-complete, and $A \pr B$, we can conclude that $B$ is $\NP$-hard. Note: You don't have to provide a formal proof. You can briefly explain in words both implications and why they hold. 
\end{enumerate}
\end{itemize}
The second bullet point above is prove that \textit{problem B} is $\NP$-hard which combined with the first bullet point yields the $\NP$-complete proof. \
\begin{itemize}
\end{itemize}

\newpage

\textbf{\large Problem 1}\\ {(30 points)} \\\\ 
Show that the following \texttt{Party-Invitation} problem is $\NP$-Complete. \\\\
\textbf{Input}: You are given a set of guests $G$ of length $m$ to be invited to a party, and a set of $k$ individual family sets $F_1, F_2, ..., F_k$ such that they are non-disjoint and $F_i \subseteq G$.\\
\textbf{Output}: A set $Q$ of guests such that not all members of any one family $F_i$ are invited to the party.\\\\
\textit{Note}: Write the reduction in the format described on the front page of this document. \\

\begin{solution}
    [Write your answer here] \vspace{12pt} % delete this line and replace with your answer/work
\end{solution}
\newpage

\textbf{\large Problem 2}\\ {(30 points)} \\\\ 
Show that the following \texttt{Average-Sum} problem is $\NP$-Complete. \\\\
\textbf{Input}: a set $S$ of natural numbers.\\
\textbf{Output}: A subset $S' \subset S$ such that the sum of the elements in the subset $S'$ equals the average of elements in the set $S$ i.e.
 \[ \sum_{s' \in S'} s' = \frac{1}{|S|} \sum_{s \in S} s \]
\\\\  
\textit{Note}: Write the reduction in the format described on the front page of this document. \\

\begin{solution}
    [Write your answer here] \vspace{12pt} % delete this line and replace with your answer/work
\end{solution}
\newpage

\textbf{\large Problem 3}\\ {(30 points)} \\\\ 
Show that the following \texttt{CLIQUE-\&-INDEPENDENT-SET} problem is $\NP$-Complete. \\\\
\textbf{Input}: an undirected graph $G = (V,E)$ and an integer $k$.\\
\textbf{Output}: Returns both a $k-$clique and $k-$independent set in $G$ i.e. a set $Q \subseteq V$ of $k$ vertices that are fully connected, AND another set $R \subseteq V$ of $k$ vertices with no edges between them. Sets $R$ and $Q$ might not be disjoint. 
\\\\
\textit{Note}: Write the reduction in the format described on the front page of this document. \\

\begin{solution}
    [Write your answer here] \vspace{12pt} % delete this line and replace with your answer/work
\end{solution}
\newpage

\textbf{\large Problem 4}\\ {(30 points)} \\\\ 
Show that the following \texttt{CYCLE-FREE} problem is $\NP$-Complete. \\\\
\textbf{Input}: a directed graph $G = (V,E)$ and an integer $k$.\\
\textbf{Output}: Returns a set $M \subseteq E$ of size at most $k$ such that $|M| \leq k$ and a subgraph $G' = (V, E \texttt{\textbackslash}M)$ is cycle-free. 
\\\\
\textit{Note}: Write the reduction in the format described on the front page of this document. \\

\begin{solution}
    [Write your answer here] \vspace{12pt} % delete this line and replace with your answer/work
\end{solution}
\newpage

\textbf{\large Problem 5} (\textit{\large Optional, but highly recommended for the exam})\\\\ 
\textit{Note: We won't grade this problem, but it's highly recommended that you solve the problem yourself.} \\\\ 
Show that the following \texttt{4-coloring} problem is $\NP$-Complete.
\\\\
\textbf{Input}: An undirected graph $G=(V,E)$ and an integer $k = 4$. \\
\textbf{Output}: A $4$-coloring of $G$ if exists, else output NO

\\\\
\newpage
    
\end{document}



